\section{Qualitative Analysis}\label{sec:exp:qualitative}

The tables average over hundreds of windows.
This section looks at the pulses themselves, first as input features and then as embeddings, to check that those averages correspond to a geometry a reader can inspect.
The input is the same for every model, so the feature plots are shown once.
Embeddings differ by model and by branch.

\subsection*{Input window}\label{sec:exp:qualitative:features}

\Cref{fig:exp:feat_w100_em,fig:exp:feat_w100_md} are the same $L=2000$ pulses from test scenario $1752$, evaluation index $100$ on the original overlapping dump used for the qualitative figures.
The four panels are carrier frequency, amplitude, pulse width, and a Cartesian incidence coordinate, each against window-local \gls{toa}.
The export names that last panel Theta. It is a direction component after \cref{eq:method:aoa_scale}, not a raw azimuth.
Color is the only change between the two figures: \cref{fig:exp:feat_w100_em} uses the recording-local emitter label, \cref{fig:exp:feat_w100_md} the global mode label.
Vanilla, RoPE-TOA, RoPE-az, Bias, and RoPE-TOA+Bias score $\ARI=1$ on both branches of this window, with two emitters and four modes.

The emitter coloring shows one platform with a scanning amplitude (the lower, scalloped track) and one platform whose amplitude sits on several flat levels.
Frequency and pulse width already separate the scanner from the other platform, but they do not separate the flat tracks from each other: two of those tracks share a frequency band near $1.45$ and the same narrow pulse width.
Incidence does separate them.
The mode coloring splits those flat tracks into different colors, which is the intended distinction.
The scanner stays one color in both plots, so in this window that emitter uses a single mode.

\begin{figure}[tbp]
  \centering
  \expfig{feat_w100_em.png}
  \caption{Test scenario $1752$ (evaluation index $100$ on the original overlapping dump), colored by ground-truth emitter. Two emitters are present.}
  \label{fig:exp:feat_w100_em}
\end{figure}

\begin{figure}[tbp]
  \centering
  \expfig{feat_w100_md.png}
  \caption{The same pulses as \cref{fig:exp:feat_w100_em}, colored by ground-truth mode. Four modes are present.}
  \label{fig:exp:feat_w100_md}
\end{figure}

Window $100$ on that dump is a light mixture.
\Cref{fig:exp:feat_maxmode} is a crowded counterpart: an overlapping evaluation window with $17$ distinct mode labels on four emitters, equal to the non-overlapping maximum used as $M$ in \cref{sec:exp:datasets}.
Frequency and pulse width collapse to a handful of horizontal bands that several modes share.
Amplitude is no longer a set of flats. Several tracks scan or hop.
Incidence still stacks platforms on separate levels, but the levels are closer and more numerous than in \cref{fig:exp:feat_w100_em}.
A method that clustered on one \gls{pdw} coordinate would not recover this window.
The dual-branch encoder still has to do so from the same four channels.

\begin{figure}[tbp]
  \centering
  \expfig{feat_maxmode.png}
  \caption{A crowded overlapping test window (evaluation index $4126$), colored by ground-truth mode.}
  \label{fig:exp:feat_maxmode}
\end{figure}

\subsection*{Embeddings}\label{sec:exp:qualitative:embed}

\Cref{fig:exp:emb_w100_em,fig:exp:emb_w100_md} are \gls{umap} projections of the per-pulse embeddings on the same scenario-$1752$ window.
Each panel is one model.
Color is again the ground-truth label of that branch (two emitters, four modes).
The source plots have no legend. The two-color panels match \cref{fig:exp:feat_w100_em}, the four-color panels match \cref{fig:exp:feat_w100_md}.
The projection is used only to look at the embeddings. It is not part of decoding.

On the emitter branch the two colors occupy separate regions for every model.
Vanilla keeps one emitter as several small islands and the other as a connected, more spread component, which matches the scanner versus the multi-level platform in the features.
RoPE-TOA and Bias pull those islands tighter and leave a wider gap between the two colors.
RoPE-az remains closer to Vanilla: the two colors still split, but the layout is less compact.
RoPE-TOA+Bias also splits the two colors on this window. The orange side is a set of thin arcs, closer to RoPE-TOA than to Vanilla.
This window is not where that model fails. The splits in \cref{tab:exp:attn_em,fig:exp:kcorr_em} sit elsewhere.

On the mode branch the four colors separate in every panel.
Some colors form thin arcs rather than round blobs, which is consistent with a mode whose amplitude or timing still moves inside the window.
None of the five joint models disagrees on this window in $\ARI$, so the visual differences are layout, not a change in the discrete partition.

\begin{figure}[p]
  \centering
  \begin{subfigure}[t]{0.32\linewidth}
    \centering
    \expfig{emb_vanilla_w100_em.png}
    \caption{Vanilla}
    \label{fig:exp:emb_w100_em:vanilla}
  \end{subfigure}\hfill
  \begin{subfigure}[t]{0.32\linewidth}
    \centering
    \expfig{emb_rope_toa_w100_em.png}
    \caption{RoPE-TOA}
    \label{fig:exp:emb_w100_em:rope_toa}
  \end{subfigure}\hfill
  \begin{subfigure}[t]{0.32\linewidth}
    \centering
    \expfig{emb_rope_az_w100_em.png}
    \caption{RoPE-az}
    \label{fig:exp:emb_w100_em:rope_az}
  \end{subfigure}\\[0.4em]
  \begin{subfigure}[t]{0.32\linewidth}
    \centering
    \expfig{emb_bias_w100_em.png}
    \caption{Bias}
    \label{fig:exp:emb_w100_em:bias}
  \end{subfigure}\hspace{0.02\linewidth}
  \begin{subfigure}[t]{0.32\linewidth}
    \centering
    \expfig{emb_combined_w100_em.png}
    \caption{RoPE-TOA+Bias}
    \label{fig:exp:emb_w100_em:combined}
  \end{subfigure}
  \caption{UMAP of emitter-branch embeddings on test scenario $1752$, colored by ground-truth emitter.}
  \label{fig:exp:emb_w100_em}
\end{figure}

\begin{figure}[p]
  \centering
  \begin{subfigure}[t]{0.32\linewidth}
    \centering
    \expfig{emb_vanilla_w100_md.png}
    \caption{Vanilla}
    \label{fig:exp:emb_w100_md:vanilla}
  \end{subfigure}\hfill
  \begin{subfigure}[t]{0.32\linewidth}
    \centering
    \expfig{emb_rope_toa_w100_md.png}
    \caption{RoPE-TOA}
    \label{fig:exp:emb_w100_md:rope_toa}
  \end{subfigure}\hfill
  \begin{subfigure}[t]{0.32\linewidth}
    \centering
    \expfig{emb_rope_az_w100_md.png}
    \caption{RoPE-az}
    \label{fig:exp:emb_w100_md:rope_az}
  \end{subfigure}\\[0.4em]
  \begin{subfigure}[t]{0.32\linewidth}
    \centering
    \expfig{emb_bias_w100_md.png}
    \caption{Bias}
    \label{fig:exp:emb_w100_md:bias}
  \end{subfigure}\hspace{0.02\linewidth}
  \begin{subfigure}[t]{0.32\linewidth}
    \centering
    \expfig{emb_combined_w100_md.png}
    \caption{RoPE-TOA+Bias}
    \label{fig:exp:emb_w100_md:combined}
  \end{subfigure}
  \caption{UMAP of mode-branch embeddings on test scenario $1752$, colored by ground-truth mode.}
  \label{fig:exp:emb_w100_md}
\end{figure}
